\documentclass[11pt,a4paper]{article}

% ----------------------------------
% Packages
% ----------------------------------
\usepackage[utf8]{inputenc}
\usepackage[T1]{fontenc}
\usepackage[a4paper,top=2.1cm,bottom=2.1cm,left=2.25cm,right=2.25cm]{geometry}
\usepackage{xcolor}
\usepackage{hyperref}
\usepackage{enumitem}
\usepackage{booktabs}
\usepackage{array}
\usepackage{titlesec}
\usepackage{fancyhdr}
\usepackage{listings}
\usepackage{graphicx}
\usepackage{float}

% ----------------------------------
% Document information
% ----------------------------------
\newcommand{\DocumentType}{Operating Systems Project Report}
\newcommand{\ProjectTitle}{Introduction to Linux Kernel Modules}
\newcommand{\CourseName}{CSE 323: Operating Systems}
\newcommand{\UniversityName}{Zagazig University}
\newcommand{\FacultyName}{Faculty of Engineering}
\newcommand{\DepartmentName}{Computer and Systems Engineering Department}
\newcommand{\AuthorName}{Youssef Reda}
\newcommand{\AcademicYear}{2026}

% ----------------------------------
% Colors and hyperlinks
% ----------------------------------
\definecolor{primary}{RGB}{25,55,95}
\definecolor{secondary}{RGB}{75,95,120}
\definecolor{linkblue}{RGB}{35,85,155}
\definecolor{codebg}{RGB}{247,248,250}
\definecolor{codeframe}{RGB}{210,216,224}

% ----------------------------------
% Code listings
% ----------------------------------
\lstset{
    basicstyle=\ttfamily\small,
    backgroundcolor=\color{codebg},
    frame=single,
    rulecolor=\color{codeframe},
    framesep=7pt,
    xleftmargin=4pt,
    xrightmargin=4pt,
    aboveskip=8pt,
    belowskip=8pt,
    breaklines=true,
    breakatwhitespace=false,
    columns=fullflexible,
    keepspaces=true,
    showstringspaces=false,
    tabsize=4,
    captionpos=b
}

\lstdefinestyle{terminal}{
    basicstyle=\ttfamily\small,
    backgroundcolor=\color{codebg},
    frame=single,
    rulecolor=\color{codeframe},
    framesep=7pt,
    breaklines=true,
    columns=fullflexible,
    keepspaces=true
}

\hypersetup{
    colorlinks=true,
    linkcolor=primary,
    urlcolor=linkblue,
    citecolor=primary,
    pdftitle={\ProjectTitle\space - \DocumentType},
    pdfauthor={\AuthorName}
}

\urlstyle{same}

% ----------------------------------
% Paragraphs, lists, and headings
% ----------------------------------
\setlength{\parindent}{0pt}
\setlength{\parskip}{4pt}

\setlist[itemize]{
    leftmargin=1.5em,
    itemsep=2pt,
    topsep=3pt
}

\setlist[enumerate]{
    leftmargin=1.8em,
    itemsep=2pt,
    topsep=3pt
}

\titleformat{\section}
    {\Large\bfseries\color{primary}}{}{0pt}{}

\titleformat{\subsection}
    {\large\bfseries\color{secondary}}{}{0pt}{}

\titleformat{\subsubsection}
    {\normalsize\bfseries}{}{0pt}{}

\titlespacing*{\section}{0pt}{16pt}{6pt}
\titlespacing*{\subsection}{0pt}{11pt}{4pt}
\titlespacing*{\subsubsection}{0pt}{8pt}{3pt}

% Unnumbered headings included in the Table of Contents
\newcommand{\nsection}[1]{%
    \section*{#1}%
    \phantomsection%
    \addcontentsline{toc}{section}{#1}%
}

\newcommand{\nsubsection}[1]{%
    \subsection*{#1}%
    \phantomsection%
    \addcontentsline{toc}{subsection}{#1}%
}

\newcommand{\nsubsubsection}[1]{%
    \subsubsection*{#1}%
    \phantomsection%
    \addcontentsline{toc}{subsubsection}{#1}%
}

% ----------------------------------
% Screenshot helper
% ----------------------------------
% Usage:
% \screenshot{assets/image-name.png}{Caption}
\newcommand{\screenshot}[2]{%
    \begin{figure}[H]
        \centering
        \includegraphics[width=0.92\textwidth,height=0.48\textheight,keepaspectratio]{#1}
        \caption{#2}
    \end{figure}
}

% ----------------------------------
% Header and footer
% ----------------------------------
\pagestyle{fancy}
\fancyhf{}
\fancyhead[L]{\small\textcolor{secondary}{Operating Systems Project}}
\fancyhead[R]{\small\textit{Linux Kernel Modules}}
\fancyfoot[C]{\thepage}

\renewcommand{\headrulewidth}{0.4pt}
\renewcommand{\footrulewidth}{0pt}
\setlength{\headheight}{14pt}

\begin{document}

% ----------------------------------
% Cover page
% ----------------------------------
\begin{titlepage}
    \thispagestyle{empty}
    \centering

    \vspace*{1.2cm}

    {\Large\bfseries\color{primary}\UniversityName\par}
    \vspace{0.25cm}
    {\large\FacultyName\par}
    {\large\DepartmentName\par}

    \vspace{1.4cm}

    {\Large\bfseries\color{secondary}\CourseName\par}
    \vspace{0.8cm}

    {\Huge\bfseries\color{primary}\ProjectTitle\par}

    \vspace{0.6cm}
    {\large Programming Project \#1\par}

    \vspace{2cm}

    \begin{center}
        \renewcommand{\arraystretch}{1.5}
        \begin{tabular}{|p{6.5cm}|p{3.5cm}|}
            \hline
            \textbf{Name}                & \textbf{Code}  \\
            \hline
            Youssef Reda Elsayed Mouawad & 20812023100544 \\
            \hline
        \end{tabular}
    \end{center}

    \vfill

    \begin{flushleft}
        \centering
        \textbf{Academic Year}\\
        \vspace{0.4cm}
        \AcademicYear
    \end{flushleft}
\end{titlepage}

% ==================================
% Table of contents
% ==================================
\pagenumbering{roman}
\tableofcontents
\clearpage
\pagenumbering{arabic}

% ==================================
% 1. Introduction
% ==================================
\nsection{Introduction to Linux Kernel Modules}

\vspace{0.5cm}
\nsubsection{Project Overview}

Linux kernel modules are pieces of code that can be loaded into
and removed from a running Linux kernel without rebuilding or
restarting the entire operating system. They provide a practical
way to extend kernel functionality and interact with kernel
facilities.

The objective of this project is to gain practical experience
with Linux kernel module development, compilation, loading,
removal, and communication with user space through the
\texttt{/proc} virtual filesystem.

The project implements three modules:

\begin{itemize}
    \item \textbf{proc-hello:} Creates a \texttt{/proc/hello}
          entry that returns a greeting message.
    \item \textbf{proc-jiffies:} Creates a \texttt{/proc/jiffies}
          entry that reports the current value of the kernel's
          \texttt{jiffies} counter.
    \item \textbf{proc-seconds:} Creates a \texttt{/proc/seconds}
          entry that reports the elapsed time in seconds since the
          module was loaded.
\end{itemize}

\vspace{0.5cm}
\nsubsection{Module Initialization and Cleanup}

A kernel module typically defines two functions: an initialization
function and a cleanup function.

\begin{itemize}
    \item \textbf{Initialization:} Runs when the module is loaded.
          It performs the required setup and returns zero on success.
    \item \textbf{Cleanup:} Runs when the module is removed.
          It releases resources allocated by the module, such as entries
          created in the \texttt{/proc} filesystem.
\end{itemize}

The \texttt{module\_init()} and \texttt{module\_exit()} macros
register these functions with the kernel. Kernel messages can be
written using \texttt{pr\_info()} or \texttt{printk()} and inspected
with the \texttt{dmesg} command.

\vspace{0.5cm}
\nsubsection{Compilation Using Make}

External kernel modules are commonly compiled using the kernel
build system, known as Kbuild. Each module directory contains
a Makefile that specifies the module's object file and invokes
the build system.

The following commands illustrate the build process:

\begin{lstlisting}[style=terminal]
make
make clean
\end{lstlisting}

The \texttt{make} command compiles the module against the build
files for the running kernel. The \texttt{make clean} command
removes generated build artifacts.

\newpage
\nsubsection{Loading, Inspecting, and Removing Modules}

The following commands are used to manage a compiled module.
Run them from the directory containing the corresponding
\texttt{.ko} file.

\begin{lstlisting}[style=terminal]
sudo insmod proc_hello.ko
lsmod
dmesg | tail
sudo rmmod proc_hello
\end{lstlisting}

\begin{itemize}
    \item \texttt{insmod} inserts a compiled module into the kernel.
    \item \texttt{lsmod} lists currently loaded kernel modules.
    \item \texttt{dmesg} displays messages from the kernel log buffer.
    \item \texttt{rmmod} removes a loaded module.
\end{itemize}

The module name passed to \texttt{rmmod} is the module's name,
not the filename with the \texttt{.ko} extension. Administrative
privileges are generally required to load and remove modules.

\vspace{0.5cm}
\nsubsection{The /proc Filesystem}

The \texttt{/proc} filesystem is a virtual filesystem that provides
an interface for accessing selected kernel and process information.
Its entries are generated dynamically and do not behave like
ordinary files stored on disk.

A kernel module can create a \texttt{/proc} entry and provide a
read callback that is invoked when a user-space program reads it.
For example:

\begin{lstlisting}[style=terminal]
cat /proc/hello
\end{lstlisting}

In the implemented modules, the callback prepares a message in
kernel memory and uses \texttt{copy\_to\_user()} to transfer the
data to a user-space buffer. The module's cleanup function removes
the corresponding \texttt{/proc} entry.

\vspace{0.5cm}
\nsubsection{Jiffies and the HZ Constant}

The kernel maintains a counter named \texttt{jiffies}, which
represents timer ticks since system boot. The \texttt{HZ} constant
defines the configured number of timer ticks per second.

The elapsed time in seconds can therefore be calculated using:

\begin{lstlisting}
elapsed_seconds = (jiffies - start_jiffies) / HZ;
\end{lstlisting}

Here, \texttt{start\_jiffies} stores the counter value captured
when the module is loaded. Subtracting it from the current
counter gives the number of ticks elapsed during the module's
lifetime. Dividing by \texttt{HZ} converts the result to seconds.

% ==================================
% 2. proc-hello
% ==================================
\newpage
\nsection{Module 1: proc-hello}

\vspace{0.5cm}
\nsubsection{Objective}

The purpose of this module is to create a \texttt{/proc/hello}
entry that returns a fixed greeting message whenever it is read.
This module demonstrates the basic interaction between a kernel
module and a user-space application through the \texttt{/proc}
filesystem.

\vspace{0.5cm}
\nsubsection{Implementation}

During initialization, the module creates the \texttt{/proc/hello}
entry and registers a read callback using \texttt{proc\_create()}
and \texttt{struct proc\_ops}. If creation fails, initialization
returns an error.

When the entry is read, the callback prepares the greeting message
and copies it to the user-space buffer using \texttt{copy\_to\_user()}.
The file position is updated to indicate that the message has
been consumed. During cleanup, the entry is removed.

\vspace{0.5cm}
\nsubsection{Compilation and Execution}

Run the following commands from the \texttt{modules/proc-hello}
directory:

\begin{lstlisting}[style=terminal]
make
sudo insmod proc_hello.ko
cat /proc/hello
sudo rmmod proc_hello
\end{lstlisting}

The expected output of the read operation is:

\begin{lstlisting}[style=terminal]
Hello from the Linux Kernel via /proc!
\end{lstlisting}

The module can also be inspected through the kernel log:

\begin{lstlisting}[style=terminal]
dmesg | tail
\end{lstlisting}

\newpage
\nsubsection{Testing and Results}

The module was built and loaded into the running kernel. Reading
\texttt{/proc/hello} returned the expected greeting message.
After removing the module, the corresponding \texttt{/proc} entry
was removed as part of the cleanup process.

\vspace{1cm}

\screenshot{../assets/proc-hello.png}
{Testing the proc-hello module and reading its /proc entry.}

% ==================================
% 3. proc-jiffies
% ==================================
\newpage
\nsection{Module 2: proc-jiffies}

\vspace{0.5cm}
\nsubsection{Objective}

The purpose of this module is to create a \texttt{/proc/jiffies}
entry that reports the current value of the kernel's
\texttt{jiffies} counter when the entry is read.

\vspace{0.5cm}
\nsubsection{Implementation}

The module uses the kernel's jiffies-related definitions to access
the current counter value. Its initialization function creates
the \texttt{/proc/jiffies} entry and registers the read callback.

When a user-space application reads the entry, the callback
formats the current counter value as text and transfers it to
user space. The cleanup function removes the entry when the
module is unloaded.

The value of \texttt{jiffies} increases as timer ticks occur.
Consequently, reading the entry at different times may return
different values.

\vspace{0.5cm}
\nsubsection{Compilation and Execution}

Run the following commands from the
\texttt{modules/proc-jiffies} directory:

\begin{lstlisting}[style=terminal]
make
sudo insmod proc_jiffies.ko
cat /proc/jiffies
cat /proc/jiffies
sudo rmmod proc_jiffies
\end{lstlisting}

The output is a numeric value representing the current
\texttt{jiffies} counter. The exact value depends on the system's
uptime and timer activity, so it should not be hard-coded in
the report.

\newpage
\nsubsection{Testing and Results}

The module was compiled and loaded successfully. Reading
\texttt{/proc/jiffies} returned the current counter value, allowing
the kernel's timer-tick counter to be observed from user space.
The entry was designed to be removed when the module is unloaded.

\vspace{1cm}

\screenshot{../assets/proc-jiffies.png}
{Testing the proc-jiffies module and displaying the current counter.}

% ==================================
% 4. proc-seconds
% ==================================
\newpage
\nsection{Module 3: proc-seconds}

\vspace{0.5cm}
\nsubsection{Objective}

The purpose of this module is to create a \texttt{/proc/seconds}
entry that reports the number of seconds elapsed since the
module was loaded.

\vspace{0.5cm}
\nsubsection{Implementation}

The module stores the value of \texttt{jiffies} at initialization
in a variable named \texttt{start\_jiffies}. Whenever the
\texttt{/proc/seconds} entry is read, the callback calculates
the difference between the current counter and the stored
starting value.

The elapsed time is calculated using:

\begin{lstlisting}
elapsed_seconds = (jiffies - start_jiffies) / HZ;
\end{lstlisting}

The subtraction gives the number of timer ticks since the module
was loaded, and division by \texttt{HZ} converts that difference
to seconds. The result is formatted as text and returned to
user space. The module removes the \texttt{/proc/seconds} entry
during cleanup.

\vspace{0.5cm}
\nsubsection{Compilation and Execution}

Run the following commands from the
\texttt{modules/proc-seconds} directory:

\begin{lstlisting}[style=terminal]
make
sudo insmod proc_seconds.ko
cat /proc/seconds
sleep 2
cat /proc/seconds
sudo rmmod proc_seconds
\end{lstlisting}

The reported value should generally increase as time passes.
Because the calculation uses integer division, the value changes
in whole seconds rather than continuously.

\newpage
\nsubsection{Testing and Results}

The module was compiled and loaded successfully. Reading
\texttt{/proc/seconds} reported the elapsed time since loading.
Reading it again after a delay demonstrated the change in elapsed
time. The module's cleanup function removes its \texttt{/proc}
entry when the module is unloaded.

\vspace{1cm}

\screenshot{../assets/proc-seconds.png}
{Testing the proc-seconds module and observing elapsed time.}

% ==================================
% 5. Conclusion
% ==================================
\newpage
\nsection{Conclusion}

\vspace{0.5cm}

This project provided practical experience with developing and
managing Linux kernel modules. The three modules demonstrated
how to register initialization and cleanup functions, compile
kernel code using Kbuild, load and remove modules, and inspect
kernel messages.

The project also demonstrated how a module can expose information
to user space through the \texttt{/proc} filesystem. The
\texttt{proc-hello} module returned a fixed message,
\texttt{proc-jiffies} exposed the current timer-tick counter, and
\texttt{proc-seconds} used the difference between two counter
values to calculate elapsed time.

Overall, the project reinforced the importance of correct resource
management and careful testing when working with kernel code,
which executes within the operating system's privileged environment.

% ==================================
% References
% ==================================
\newpage
\nsection{References}

\begin{enumerate}
    \item Abraham Silberschatz, Peter B. Galvin, and Greg Gagne,
          \textit{Operating System Concepts}, 10th Edition,
          Chapter 2, Programming Project: Introduction to Linux Kernel Modules.

    \item Linux Kernel Documentation,
          \href{https://docs.kernel.org/}{The Linux Kernel Documentation}.
\end{enumerate}

\end{document}
